% !PRINTED 676
\weberpage{676}
%<<CONT
and thus obtain from (12), if $\lambda_s', \lambda_s'', \dots$ are the conjugate logarithms of $\eta', \eta'', \dots$:
\weq{14}{
\begin{aligned}
a &< \Sigma g_s \lambda_s < a'\\
a' &< \Sigma g_s \lambda_s' < a''\\
a'' &< \Sigma g_s \lambda_s'' < a'''\\
&\cdots\cdots\cdots\cdots,
\end{aligned}
}
and the sequence of these inequalities can be continued without limit.

All these numbers $\eta, \eta' \eta'', \dots$ have moreover, by (6) and (7), the property that their absolute norms $N, N', N'', \dots$ remain below a definite finite limit $e^{nk}$.

% [sic] The source omits the comma between \eta' and \eta'' in the preceding list.

The number of possible values of $N, N', N'', \dots$ is therefore finite, namely certainly not greater than the greatest integer contained in $e^{nk}$. Likewise the number of all possible residues of the numbers $x_1, x_2, \dots, x_n$ [in (4)] with respect to all the moduli $N, N', N'', \dots$ is only a finite one, and from this it follows that, if we continue the sequence of the numbers (13) only far enough, there must occur in the sequence \gesperrt{two distinct numbers $\eta^{(h)}, \eta^{(k)}$ in which $N^{(h)} = N^{(k)}$ and the residues of $x_1, x_2, \dots, x_n$ with respect to the modulus $N^{(h)}$ are exactly the same.}

From (14), however, it follows, if $h > k$ is assumed:
\weq{15}{\Sigma g_s (\lambda_s^{(h)} - \lambda_s^{(k)}) > 0.}
If $N$ is the absolute norm of these two numbers $\eta^{(h)}, \eta^{(k)}$, then, because of the agreement of the residues of the $x$, the difference $\eta^{(h)} - \eta^{(k)}$ is divisible by $N$. On the other hand, by \S.\,138\,(13) the number $N$ is divisible by $\eta^{(h)}$, and consequently $\eta^{(h)} - \eta^{(k)}$ is also divisible by $\eta^{(h)}$. From this it results that $\eta^{(k)}$ is also divisible by $\eta^{(h)}$, and likewise $\eta^{(h)}$ by $\eta^{(k)}$. Both numbers are therefore associated, and if we set
\[
\frac{\eta^{(h)}}{\eta^{(k)}} = \varepsilon
\]
then $\varepsilon$ is a \gesperrt{unit} of the field $\Omega$.

If we further set, denoting by $|\varepsilon_s|$ the absolute value of $\varepsilon_s$,
\weq{16}{\delta_s \log|\varepsilon_s| = l_s,}
then from (15) there results:
\weq{17}{g_1 l_1 + g_2 l_2 + \cdots + g_\nu l_\nu > 0.}
